\section{Worked examples}
\label{sec:worked-examples}

\subsection{A complete finite-sum trace}

Let \(f(p,i)\) be a fixed primitive-recursive integer-valued function and let
\(n(p)\) be a bound.  Define
\[
 S(p,0)=0,
 \qquad
 S(p,i+1)=S(p,i)+f(p,i).
\]
The trace state is the current signed-integer sum, the trace bound is \(n(p)\),
and the invariant at stage \(i\) is the equation
\[
 s=S(p,i).
\]
The base follows from the zero equation for \(S\).  The local step rewrites
both the trace transition and \(S(p,i+1)\) to addition of \(f(p,i)\).  The
final implication identifies the output state with \(S(p,n(p))\).

Although elementary, this is the right first Lean milestone.  It exercises
\code{PR.prec}, \code{PRA.precZero}, \code{PRA.precSucc}, equality transport,
\code{PRA.qfInd}, specialization at a primitive-recursive bound, implication
discharge, and universal introduction.  Once this proof is a reusable
\code{traceReflection} instance, finite products, Horner evaluation, rational
error accumulation, and path generation differ only in their local invariant.

\subsection{Relational square-root approximations}

Consider \(\sqrt2\), but deliberately do not define it as a completed real.
At level \(k\), let the grid denominator be
\(D_k=2^{k+4}\).  A candidate is an integer
\(a\) with \(D_k\le a\le2D_k\), interpreted as
\(q_{k,a}=a/D_k\).  Declare it good when
\begin{equation}
 |q_{k,a}^2-2|\le 2^{-k}.
 \label{eq:sqrt-good}
\end{equation}
Declare adjacent candidates compatible when
\begin{equation}
 |q_{k,a}-q_{k+1,b}|\le 3\cdot2^{-k}.
 \label{eq:sqrt-compat}
\end{equation}
All tests are exact integer inequalities after clearing powers of two.

There are generally several good grid points.  No exact decision about which
one is ``the'' approximation is needed.  Nonemptiness follows from a bounded
integer-square-root search: choose the least \(a\) with \(a^2\ge2D_k^2\), and
compare \(a/D_k\) with the preceding grid point.  The mesh is fine enough to
make one of these points satisfy \eqref{eq:sqrt-good}.

If \(q\) and \(r\) are two good candidates at level \(k\), then
\(q,r\in[1,2]\), and
\[
 |q-r|=\frac{|q^2-r^2|}{q+r}
 \le \frac{|q^2-2|+|r^2-2|}{2}
 \le 2^{-k}.
\]
Thus the viable cloud has diameter at most \(2^{-k}\), independently of which
finite path produced its candidates.  Compatibility gives an edge budget
\(3\cdot2^{-k}\).  A slightly larger precision shift makes the geometric tail
and the horizontal diameter smaller than the regular-name target.  The
shrinking-viability theorem therefore turns the least full-viable grid point
into a primitive-recursive Cauchy name.

The executable companion enumerates these finite grids, computes full-viable
clouds by dynamic programming, and checks the resulting bounds for sample
horizons.  This is regression evidence for the definitions, not a replacement
for the universal iPRA argument above.

\subsection{Roots at a collision}

Consider
\[
 P_t(Y)=Y^2-t^2.
\]
For positive \(t\), the two roots can be labeled by sign.  At \(t=0\), they
collide.  A uniform ordered pair of roots whose first component is always the
negative root and second is the positive root requires an exact sign decision
on represented \(t\) if the presentation permits both signs and zero.

The relational presentation returns two root boxes, counted with
multiplicity, or one cluster of multiplicity two when the separation is below
the requested scale.  Polynomial residuals and Vieta coefficient checks are
finite.  If a certificate supplies \(|t|\ge g>0\), then approximating \(t\)
within \(g/3\) yields its sign, separates the roots by \(2g\), and recovers the
ordered labels.  Thus the same library supports:

\begin{itemize}
\item unconditional multiset/cluster output;
\item an individual approximate root relation;
\item exact ordered roots under a positive separation guard.
\end{itemize}

This pattern extends to eigenvalues, eigenvectors, algebraic branches, and
projective directions.

\subsection{Overlapping coordinate charts}

Suppose a represented point on a curve can be computed in chart \(U_0\) when
\(x<2/3\) and in chart \(U_1\) when \(x>1/3\).  The overlap
\((1/3,2/3)\) is intentional.  Approximate \(x\) within \(1/12\) and compare
the approximation with \(1/2\).  If it is below \(1/2\), use \(U_0\);
otherwise use \(U_1\).  A wrong comparison can occur only when the true point
lies in the overlap, where both formulas are certified.

The branch bit is not an invariant mathematical object, but the output point
is extensional.  This small example captures the preferred treatment of
pivots, local logarithms, rational parametrizations, and analytic charts.

\subsection{Bounded rational points}

Let \(F_p(X,Y)\in\Z[X,Y]\) be a fixed primitive-recursive coefficient family.
A normalized rational pair of height at most \(H\) is coded by integers
\((a,b,c)\) with \(c>0\), bounded absolute values, and the selected gcd/sign
normalization.  The equation
\(F_p(a/c,b/c)=0\) becomes an exact integer equation after multiplying by a
fixed power of \(c\).

The finite enumerator loops over the bounded box, retains normalized solutions,
and removes duplicates by code equality.  Its invariant says that the current
list contains exactly the solutions among the first \(j\) candidate codes.
The trace compiler proves soundness and completeness below \(H\).

A separate height certificate
\[
 F_p(x,y)=0\to h(x,y)\le H
\]
turns this bounded list into the complete list.  Without the certificate, the
same executable search is only a list below a cutoff.  This is the precise
status distinction needed in the Faltings and Siegel settings.

\subsection{A finite-net minimum}

Let \(f:[0,1]^d\to\R\) have a supplied rational modulus.  At requested
precision \(k\), choose a rational grid fine enough that variation inside a
cell is at most \(2^{-k-2}\).  Evaluate \(f\) at the grid points to error
\(2^{-k-2}\), and return the least point with minimum approximate value.
Finite search proves that its displayed value is within \(2^{-k}\) of the
infimum.

If there may be several minimizers, the point itself remains a relation of
near-minimizers.  If a strong-convexity certificate gives
\[
 f(x)\ge f(x_*)+\lambda\|x-x_*\|^2,
\]
then every near-minimizer lies in a rational ball of radius obtained from its
value error.  These balls shrink, so the least full-viable grid point is a
coherent name for \(x_*\).  The value theorem and point theorem are thereby
cleanly separated.

\subsection{End-to-end proof flow}

A mature packet follows the flow
\[
\boxed{\text{codes and guards}}
\longrightarrow
\boxed{\text{finite observations}}
\longrightarrow
\boxed{\text{trace/certificate}}
\longrightarrow
\boxed{\Prf\,\PRAAx}
\longrightarrow
\boxed{\NPrf}.
\]
If outputs are non-unique, insert
\[
\boxed{\text{finite relation}}
\longrightarrow
\boxed{\text{viability}}
\longrightarrow
\boxed{\text{shrinking, if justified}}
\]
between finite observations and the final represented object.  If a deep
certificate producer is unknown, the first four boxes can still be completed
as a conditional theorem whose missing input is explicit.
